\section*{अध्याय \ref{ch:b3:inorganic-materials} --- अकार्बनिक पदार्थ और नैनोपदार्थ}

\begin{solution}{exo:b3:inorganic-materials:1}
$x^2 \propto t$: तीन गुना मोटाई में नौ गुना समय लगता है, \qty{90}{h}।
\end{solution}

\begin{solution}{exo:b3:inorganic-materials:2}
जल-अपघटन: $\ce{Si(OC2H5)4 + 4H2O -> Si(OH)4 + 4C2H5OH}$।\\
संघनन: $\ce{2Si(OH)4 -> Si2O(OH)6 + H2O}$।\\
समग्र: $\ce{Si(OC2H5)4 + 2H2O -> SiO2 + 4C2H5OH}$।
\end{solution}

\begin{solution}{exo:b3:inorganic-materials:3}
निर्माता: $\ce{SiO2}$, $\ce{B2O3}$। रूपांतरक: $\ce{Na2O}$, $\ce{CaO}$। $\ce{Na2O}$ का प्रत्येक ऑक्साइड आयन एक Si--O--Si
सेतु को दो Si--O$^-$ सिरों में तोड़ता है, प्रत्येक एक $\ce{Na+}$ के साथ युग्मित: जाल कट जाता है, गलित
अधिक आसानी से बहता है और निम्नतर ताप पर पिघलता है।
\end{solution}

\begin{solution}{exo:b3:inorganic-materials:4}
बारह पंचभुज (हर \omterm{def:b3:inorganic-materials:carbon}{फ़ुलरीन} की तरह)। $V = 70$, $E = 3V/2 = 105$ आबंध, $F = 2 - V + E = 37$
फलक, अतः $37 - 12 = 25$ षट्भुज।
\end{solution}

\begin{solution}{exo:b3:inorganic-materials:5}
$\sqrt2\,(r_B + r_O) = 1.4142 \times \qty{200.5}{pm} = \qty{283.5}{pm}$। $\ce{SrTiO3}$: $284/283.5 = 1.00$, आदर्श घनीय पेरोव्स्काइट।
$\ce{BaTiO3}$: $301/283.5 = 1.06$, टाइटेनियम अपने अष्टफलक के लिए बहुत छोटा, केंद्र से हटा हुआ: \omterm{def:b3:inorganic-materials:ferroelectric}{लौहविद्युत}।
$\ce{CaTiO3}$: $274/283.5 = 0.97$, कैल्सियम अपनी गुहा के लिए बहुत छोटा: झुके अष्टफलक, निम्नतर
सममिति।
\end{solution}

\begin{solution}{exo:b3:inorganic-materials:6}
Si/Al $= 106/86 = 1.23$। $M = 86 \times 23.0 + 86 \times 27.0 + 106 \times 28.1 + 384 \times 16.0 = \qty{13422.6}{g/mol}$; प्रति सूत्र 43 $\ce{Ca^2+}$:
$43/\qty{13422.6}{g/mol} = \qty{3.20}{mmol/g}$।
\end{solution}

\begin{solution}{exo:b3:inorganic-materials:7}
निकटतम-पड़ोसी दूरी $d = a/\sqrt2 = \qty{288.4}{pm}$। $n$ कोशों वाला गुच्छ $(2n + 1)d$ चौड़ा होता है:
$2n + 1 \approx 5/0.2884 = 17.3$, अतः $n = 8$ (\qty{4.9}{nm})। $N(8) = 2057$ परमाणु, $10 \times 64 + 2 = 642$ पृष्ठ पर:
31\,\%।
\end{solution}

\begin{solution}{exo:b3:inorganic-materials:8}
$h^2/(8\mu R^2)$, जहाँ $\mu = \qty{9.109e-32}{kg}$: \qty{2.0}{nm} पर \qty{0.94}{eV} और \qty{3.0}{nm} पर \qty{0.42}{eV}। उत्सर्जन
$1.74 + 0.94 = \qty{2.68}{eV}$ पर, \qty{463}{nm} (नीला), और $1.74 + 0.42 = \qty{2.16}{eV}$ पर, \qty{574}{nm} (पीला): छोटा बिंदु
अधिक नीला प्रकाश उत्सर्जित करता है।
\end{solution}

\begin{solution}{exo:b3:inorganic-materials:9}
$(10, 10)$: आरामकुर्सी, $n - m = 0$, धात्विक। $(10, 0)$: ज़िगज़ैग, 10, 3 का गुणज नहीं, \omterm{def:b3:solid-state:classes}{अर्धचालक}।
$(9, 0)$: ज़िगज़ैग, धात्विक। $(12, 8)$: दोनों में से कोई नहीं (किरैल), $n - m = 4$, \omterm{def:b3:solid-state:classes}{अर्धचालक}।
\end{solution}

\begin{solution}{exo:b3:inorganic-materials:10}
$M = 4 \times 65.4 + 13 \times 16.0 + 24 \times 12.0 + 12 \times 1.0 = \qty{769.6}{g/mol}$; $a^3 = (\qty{25.82e-8}{cm})^3 = \qty{1.721e-20}{cm^3}$;
$\rho = 8 \times 769.6/(\num{6.022e23} \times \num{1.721e-20}) = \qty{0.594}{g.cm^{-3}}$। एक मैदान \qty{7140}{m^2} है: एक ग्राम में
$3000/7140 = 0.42$ मैदान हैं, और कुछ ग्राम के एक चम्मच में पूरे मैदान से अधिक।
\end{solution}

\begin{solution}{exo:b3:inorganic-materials:11}
$\dfrac{10n^2 + 2}{(10n^3 + 15n^2 + 11n + 3)/3} = \dfrac{30n^2 + 6}{10n^3 + 15n^2 + 11n + 3} \to \dfrac{30n^2}{10n^3} = \dfrac3n$। $n = 8$ के लिए: यथार्थ
$642/2057 = 0.31$, जबकि $3/8 = 0.375$; सन्निकटन छोटे गुच्छों के लिए अधिक आकलन करता है,
जिनके भीतरी कोश नगण्य नहीं हैं।
\end{solution}

\begin{solution}{exo:b3:inorganic-materials:12}
$2E = 3V$, $2E = 5p + 6h + 7s$, $F = p + h + s$। ऑयलर: $V - E + F = 2$, $V = 2E/3$ के साथ: $F - E/3 = 2$, अतः
$p + h + s - (5p + 6h + 7s)/6 = 2$, अर्थात् $(p - s)/6 = 2$: $p - s = 12$। प्रत्येक सप्तभुज को एक अतिरिक्त
पंचभुज चाहिए।
\end{solution}

\begin{solution}{pb:b3:inorganic-materials:1}
\textbf{1.} $\mathrm{Na_{12}Al_{12}Si_{12}O_{48}}$: $12 \times 23.0 + 12 \times 27.0 + 12 \times 28.1 + 48 \times 16.0 = \qty{1705.2}{g/mol}$।

\textbf{2.} $1705.2 + 27 \times 18.0 = \qty{2191.2}{g/mol}$।

\textbf{3.} Si/Al $= 12/12 = 1$।

\textbf{4.} कोई Al--O--Al कड़ी नहीं; Si/Al $= 1$ के साथ प्रत्येक Si के चार Al पड़ोसी हैं और प्रत्येक
Al के चार Si: कठोर एकांतरण।

\textbf{5.} $-12$, प्रति ऐलुमिनियम एक।

\textbf{6.} $486.0/2191.2 = 22.2$\,\%।

\textbf{7.} $\ce{2Na+}(\text{ज़िओलाइट}) + \ce{Ca^2+}(\text{aq}) \rightleftharpoons \ce{Ca^2+}(\text{ज़िओलाइट}) + \ce{2Na+}(\text{aq})$।

\textbf{8.} छह (बारह आवेश)।

\textbf{9.} $6/\qty{1705.2}{g/mol} = \qty{3.52}{mmol/g}$।

\textbf{10.} $6/\qty{2191.2}{g/mol} = \qty{2.74}{mmol/g}$।

\textbf{11.} प्रति ग्राम $\qty{3.52}{mmol/g} \times \qty{100.1}{g/mol} = \qty{352}{mg}$ $\ce{CaCO3}$।

\textbf{12.} $\qty{2.0}{mmol/L} \times \qty{15}{L} = \qty{30}{mmol}$।

\textbf{13.} $\qty{30}{mmol}/\qty{2.74}{mmol/g} = \qty{11}{g}$।

\textbf{14.} $\qty{10.9}{g}/0.60 = \qty{18}{g}$।

\textbf{15.} $\qty{60}{mmol}$ $\ce{Na+}$, $\qty{0.060}{mol} \times \qty{23.0}{g/mol} = \qty{1.4}{g}$।

\textbf{16.} अपशिष्ट जल में फ़ॉस्फ़ेट नदियों और झीलों में शैवाल का पोषण करते हैं; शैवाल-प्रस्फुटन
और उनका क्षय घुली हुई ऑक्सीजन को समाप्त कर देते हैं (सुपोषण)।

\textbf{17.} छोटा $\ce{Mg^2+}$ अपने जल अणुओं को अधिक दृढ़ता से थामे रखता है; खिड़की से गुज़रकर स्थलों तक पहुँचने के लिए
उसे उन्हें छोड़ना पड़ता है, जो धुलाई तापों पर धीमा है।

\textbf{18.} नग्न आयन लगभग \qty{0.20}{nm} चौड़ा है, खिड़की का आधा।

\textbf{19.} जलयोजित आयन खिड़की से बड़ा है: गुज़रने के लिए उसे अपने
जल कोश का कुछ भाग छोड़ना पड़ता है, और ढाँचे के ऑक्सीजन जल का
स्थान ले लेते हैं।

\textbf{20.} ढाँचा ऋणावेशित है और केवल धनायनों का विनिमय करता है; लंबी शृंखला वाला
बड़ा ऋणायन प्रतिकर्षित होता है और खिड़की से नहीं गुज़र सकता, अतः वह
जल में ही रहता है, जहाँ वह अपना कार्य करता है।

\textbf{21.} निर्जलित \omterm{def:b3:inorganic-materials:zeolite}{ज़िओलाइट} अपने पिंजरों में जल को प्रबलता से ग्रहण करता है; केवल
खिड़की से छोटे अणु ही प्रवेश करते हैं, अतः यह अणुओं को आकार के अनुसार पृथक करता है:
आण्विक पैमाने पर एक चलनी।

\textbf{22.} बड़े $\ce{K+}$ आयन खिड़कियों में बैठकर उन्हें संकरा कर देते हैं: केवल छोटे
अणु, जैसे जल, प्रवेश पाते हैं।

\textbf{23.} पतला, रैखिक \textit{पैरा} समावयवी चैनलों में समा जाता है और तेज़ी से विसरित होता है;
चौड़े \textit{ऑर्थो} और \textit{मेटा} समावयवी धीरे विसरित होते हैं, रुकते हैं और समावयवित होते हैं।

\textbf{24.} निर्जल \omterm{def:b3:inorganic-materials:zeolite}{ज़िओलाइट} A सैद्धांतिक रूप से प्रति ग्राम \qty{3.52}{mmol} $\ce{Ca^2+}$ ग्रहण कर सकता है।
\end{solution}
